
%! Author = Omar Iskandarani


%! Date = 3/7/2026


%! Affiliation = Independent Researcher, Groningen, The Netherlands


%! License = © 2026 Omar Iskandarani. All rights reserved. This manuscript is made available for academic reading and citation only. No republication, redistribution, or derivative works are permitted without explicit written permission from the author. Contact: info@omariskandarani.com


%! ORCID = 0009‑0006‑1686‑3961


%! DOI = 10.5281/zenodo.yyy






\newcommand{\paperdoi}{10.5281/zenodo.yyy}


\newcommand{\papertitle}{\textbf{Hydrodynamics of Swirl‑Strings:}\ \Large Emergent Quantum Constants, Topological Leptons and Discrete Scale Invariance}






%=========================================


% PREAMBLE, PACKAGES AND DOCUMENT CONFIGURATION


%=========================================


\documentclass[11pt]{article}


\usepackage{amsmath,amssymb,amsfonts,bm}


\usepackage{siunitx}


\usepackage[hidelinks]{hyperref}


\usepackage[a4paper,margin=1in]{geometry}


\usepackage[T1]{fontenc}


\usepackage[utf8]{inputenc}


\usepackage{booktabs}


\usepackage{titlesec}






% swirl arrows (context‑aware) and canonical macros


\newcommand{\swirlarrow}{\mkern-2mu\scriptscriptstyle\boldsymbol{\circlearrowleft}}


\newcommand{\vswirl}{\mathbf{v}_{\mkern-2mu\scriptscriptstyle\boldsymbol{\circlearrowleft}}}


\newcommand{\SwirlClock}{S_{(t)}^{\mkern-2mu\scriptscriptstyle\boldsymbol{\circlearrowleft}}}


\newcommand{\Fmaxswirl}{F^{\max}_{\mkern-1mu\scriptscriptstyle\boldsymbol{\circlearrowleft}}}


\newcommand{\Fmax}{F^{\max}_{\mkern-1mu\scriptscriptstyle\boldsymbol{\circlearrowleft}}} 


\newcommand{\FmaxEM}{F^{\max}_{\mathrm{EM}}}


\newcommand{\FmaxG}{F_{\mathrm{G}}^{\max}}               % G‑like maximal force scale


\newcommand{\vscore}{v_{\swirlarrow}}                    % shorthand: |v_swirl| at r=r_c


\newcommand{\vnorm}{\lVert \mathbf{v}_{\mkern-2mu\scriptscriptstyle\boldsymbol{\circlearrowleft}} \rVert}  % swirl speed magnitude


\newcommand{\rhoF}{\rho_{\!f}}\newcommand{\rhof}{\rho_{\!f}}     % effective fluid density


\newcommand{\rhoE}{\rho_{\!E}}\newcommand{\rhoe}{\rho_{\!E}}                           % swirl energy density


\newcommand{\rhoM}{\rho_{\!m}}\newcommand{\rhom}{\rho_{\!m}}                           % mass‑equivalent density


\newcommand{\omegas}{\boldsymbol{\omega}_{\swirlarrow}}  % swirl vorticity


\newcommand{\Om}{\Omega_{\swirlarrow}}                   % swirl angular frequency profile


\newcommand{\rc}{r_c}                                    % string core radius (swirl string radius)


\newcommand{\rhocore}{\rho_{\text{core}}}


\newcommand{\GammaSST}{\Gamma_0}






% additional macros for constants and notation


\newcommand{\phiG}{\varphi}


\newcommand{\me}{m_e}


\newcommand{\Eeff}{\mathcal{E}_{\rm eff}}






\titleformat{\section}{\large\bfseries}{\thesection}{1em}{}






%=========================================


% Title Page Macros


%=========================================


\newcommand{\titlepageOpen}{%


    \begin{titlepage}


        \thispagestyle{empty}  \centering


        \Large \bfseries \papertitle \par \vspace{1cm}


        {\Large \itshape \textbf{Omar Iskandarani}\textsuperscript{\textbf{*}} \par} \vspace{0.5cm}


        {\today \par}  \vspace{0.5cm}


}






\newcommand{\titlepageClose}{%


        \begin{picture}(0,0)


            \put(0,-45){  % Shift left and down


                \begin{minipage}[b]{0.7\textwidth} \footnotesize


                    \renewcommand{\arraystretch}{1.0} \noindent\rule{\textwidth}{0.4pt} \\[0.5em]


                    \textsuperscript{\textbf{*}} Independent Researcher, Groningen, The Netherlands \\


                    Email: \texttt{info@omariskandarani.com} \\


                    ORCID: \texttt{\href{https://orcid.org/0009-0006-1686-3961}{0009-0006-1686-3961}} \\


                    DOI: \href{https://doi.org/\paperdoi}{\paperdoi}


                \end{minipage}


            }


        \end{picture}


    \end{titlepage}%


}






%=========================================


% DOCUMENT


%=========================================


\begin{document}


    \titlepageOpen


        \begin{abstract}


            Swirl‑String Theory (SST) models elementary particles as knotted vortex loops in an incompressible, inviscid superfluid.  Unlike point‑particle or field‑theoretic descriptions, SST attributes mass, spin and magnetic moments to the hydrodynamic properties of structured vorticity.  In this manuscript we derive, from first principles, the electron mass, the reduced Planck constant and the anomalous magnetic moment using a Golden‑parameterized Nonlinear Schrödinger (NLS) vortex core.  We show that the self‑intersection of a trefoil vortex ring ($T_{3,2}$) collapses into a Hill spherical vortex; this ``degenerate torus'' reproduces the observed electron rest mass.  We then compute the intrinsic angular momentum and recover the physical value of $\hbar$ from the circulation quantum and core density.  Acoustic shear at the vortex boundary leads naturally to the $g$‑factor anomaly $a_e\approx \alpha/2\pi$.  Finally, we demonstrate that higher‑generation leptons correspond to discrete topological states of $T_{p,2}$ torus knots with masses scaling as $\varphi^{-2k}$, giving a quantitative description of the muon and tau.  Throughout, we compare our hydrodynamic quantities with textbook definitions of the classical electron radius【364253867482672†L183-L189】, the Compton frequency【620899862539895†L181-L186】, the Bohr magneton【733406408751†L151-L162】 and the anomalous magnetic moment【733796025034506†L150-L156】.


        \end{abstract}


    \titlepageClose






%======================================================================


\section{Introduction and Axioms}


\label{sec:introduction}


Swirl‑String Theory (SST) is predicated on the existence of a universal, incompressible and inviscid fluid medium.  Mass and gravity arise not from curvature or quantized fields but from local vorticity and pressure gradients within this medium.  A particle is represented by a closed loop of vorticity (a swirl‑string), characterised by a finite core radius $\rc$ and circulation quantum $\GammaSST$.  To avoid the singularities inherent to thin‑tube approximations, SST adopts a Nonlinear Schrödinger (NLS) vortex ring with a ``healing length'' equal to the core radius.  The requirement that the pressure at the core boundary is constant (``const druk'') uniquely selects the Golden Ratio $\phiG=(1+\sqrt{5})/2$ as a parameter of the NLS solution.






The aim of this work is twofold: (i) derive emergent quantum constants ($m_e$, $\hbar$, $\mu_B$, $g-2$) from the hydrodynamics of the swirl‑string, and (ii) elucidate the discrete scale invariance that leads to the lepton mass spectrum.  Our approach is self‑contained and does not presuppose results from quantum mechanics; instead, we use textbook physical constants as calibration points.






\paragraph{Canonical constants.}  We anchor our derivations using the classical electron radius $r_e$【364253867482672†L183-L189】, defined by $r_e = e^2/(4\pi\varepsilon_0 m_e c^2)$ or equivalently $r_e = \alpha\hbar/(m_e c)$, where $e$ is the elementary charge, $m_e$ the electron mass, $\varepsilon_0$ the vacuum permittivity, $c$ the speed of light and $\alpha$ the fine‑structure constant.  In SST we identify the core radius as $\rc = r_e/2$.  The characteristic swirl speed at the core boundary is set by the Compton frequency $\omega_c = m_e c^2/\hbar$【620899862539895†L181-L186】, leading to $\vnorm = \omega_c\,\rc$.  The Planck constant $\hbar$ and Bohr magneton $\mu_B = e\hbar/(2 m_e)$【733406408751†L151-L162】 will emerge from hydrodynamic quantities.






\paragraph{Notation.}  We use $\rhocore$ for the core density, $\rhoF$ for the effective fluid density at large scales, and $\rhoe = \rhocore c^2$ for the energy density.  All numerical evaluations will use CODATA 2025 values unless otherwise stated.






\paragraph{Structure.}  Section~\ref{sec:NLS} formulates the Golden NLS vortex and solves for the degenerate torus representing the electron.  Section~\ref{sec:emergent_constants} derives $\hbar$, the Bohr magneton and the magnetic moment anomaly from hydrodynamic first principles.  Section~\ref{sec:leptons} relates higher‑generation leptons to topological torus knots.  Appendices provide extended derivations and discuss discrete tangle complexity.






%======================================================================


\section{Golden NLS Vortex and Degenerate Torus}


\label{sec:NLS}


The energy of a vortex ring of major radius $R$ and core radius $\rc$ in a compressible superfluid can be expressed in the NLS framework as


\begin{equation}


    E = \tfrac{1}{2} \rhocore \GammaSST^2 R \left[ \ln \left( \frac{8R}{\rc} \right) - A \right],


    \label{eq:NLS_energy}


\end{equation}


where $A$ is a dimensionless constant determined by the core pressure profile.  In SST the const‑druk condition requires that the pressure at the boundary is constant; thermodynamic analysis (see Appendix~\ref{app:golden}) shows that this is satisfied by setting $A=\phiG$ and the corresponding translational coefficient $B=\phiG^{-1}$.  The Golden Ratio therefore controls the self‑similar scaling of the vortex core.






Using mass–energy equivalence $E = m c^2$, the mass functional becomes


\begin{equation}


    m(R) = \frac{\rhocore \GammaSST^2 R}{2 c^2} \left[ \ln \left( \frac{8R}{\rc} \right) - \phiG \right].


    \label{eq:mass_functional}


\end{equation}


Given the canonical triad $(\rhocore,\GammaSST,\rc)$, we solve the transcendental equation $m(R)=m_e$ for the dimensionless ratio $x=R/\rc$.  Substituting $\rhocore = 3.893\times 10^{18}\,\text{kg m}^{-3}$, $\GammaSST = 9.6845\times 10^{-9}\,\text{m}^2\text{s}^{-1}$ and $\rc = 1.40897\times 10^{-15}\,\text{m}$, Eq.~\eqref{eq:mass_functional} yields


\begin{equation}


    x\,[\ln(8x) - \phiG] \approx 0.31826,\quad x \approx 0.895.


\end{equation}


Hence the major radius is $R_e \approx 0.895\,\rc$; it is smaller than the core radius, implying that the toroidal hole collapses and the vortex forms a contiguous spherical droplet.  This configuration corresponds exactly to Hill's spherical vortex【740404540107499†L143-L149】, an exact solution of the Euler equations discovered by Hill in 1894 and often used to model vortex rings.  The resulting structure is effectively ``point‑like'' at scales much larger than $\rc$, consistent with electron scattering experiments.






\subsection{Hooke closure and maximal swirl force}


The Hooke closure postulates that the maximum swirl force within the core balances a harmonic restoring force, $m_e c^2 = \tfrac{1}{2} k r_e^2$.  With $\rc = r_e/2$, the maximal swirl force is


\begin{equation}


    \Fmax = \frac{m_e c^2}{2\rc} \approx 29.05\,\text{N},


\end{equation}


and the core density is


\begin{equation}


    \rhocore = \frac{\Fmax}{\pi \rc^2 \vnorm^2},


\end{equation}


where $\vnorm = \omega_c\,\rc$ is the characteristic swirl speed at the core boundary.  These expressions are derived and numerically validated in Appendix~\ref{app:constants}.






%======================================================================


\section{Emergent Constants and Magnetic Moments}


\label{sec:emergent_constants}


Having established the degenerate torus topology, we now derive the intrinsic angular momentum and magnetic moments of the vortex droplet.






\subsection{Planck constant from hydrodynamic spin}


The spin of the droplet arises from the circulation of the core fluid.  Treating the degenerate torus as a sphere of radius $\rc$ with density $\rhocore$, the angular momentum is the product of the moment of inertia and angular velocity.  In SST this yields


\begin{equation}


    \hbar_{\rm SST} = \rhocore \rc^3 \GammaSST.


    \label{eq:hbar_sst}


\end{equation}


Substituting numerical values gives


\begin{equation}


    \hbar_{\rm SST} = (3.893\times 10^{18})\,(2.797\times 10^{-45})\,(9.684\times 10^{-9}) \approx 1.054\times 10^{-34}\,\text{J s},


\end{equation}


in excellent agreement with the physical reduced Planck constant $\hbar$【620899862539895†L181-L186】.  Quantization thus emerges from hydrodynamic stiffness and topology rather than being postulated.






\subsection{Bohr magneton and Dirac $g=2$}


The magnetic moment of a classical current loop is $\mu = I A$.  In SST the current arises from the circulation of charge around the core.  Using Eq.~\eqref{eq:hbar_sst} in the Bohr magneton definition $\mu_B = e\hbar/(2m_e)$【733406408751†L151-L162】, the SST magnetic moment is


\begin{equation}


    \mu_{\rm SST} = \frac{e\,\hbar_{\rm SST}}{2 m_e} = \frac{e}{2 m_e}\,\rhocore \rc^3 \GammaSST.


\end{equation}


Since $\hbar_{\rm SST}$ matches $\hbar$, this reproduces the Dirac prediction for a spin‑1/2 particle with $g=2$ without invoking point charges.






\subsection{Acoustic dressing and the anomalous magnetic moment}


Standard QED attributes the anomalous magnetic moment $a_e=(g-2)/2$ to radiative corrections.  In SST the anomaly arises from acoustic shear at the boundary between the spinning core and the irrotational background.  The rotational speed at the equator is $\vscore = \alpha c/2$, where $\alpha$ is the fine‑structure constant.  Shear excites Lamb acoustic modes which modify the effective circulation; a first‑order perturbation gives


\begin{equation}


    a_e \approx \frac{\alpha}{2\pi} \approx 0.0011614,


    \label{eq:ae}


\end{equation}


in agreement with the one‑loop QED result【733796025034506†L150-L156】.  Higher‑order terms correspond to nonlinear acoustic couplings and scale with powers of $\alpha$.






The complete magnetic moment including acoustic dressing is therefore


\begin{equation}


    \mu_e = \frac{e\,\hbar_{\rm SST}}{2 m_e}\left[1 + \frac{\alpha}{2\pi} + \mathcal{O}(\alpha^2)\right].


\end{equation}






%======================================================================


\section{Topological Lepton Generations}


\label{sec:leptons}


SST predicts that heavier leptons correspond to more complex torus knot topologies.  A $(p,2)$ torus knot has $p$ meridional crossings and two longitudinal windings.  The vortex energy scales linearly with $p$ and exponentially with discrete screening layers indexed by $k$; thus the mass of a lepton with knot number $p$ and layer index $k$ is


\begin{equation}


    m(p,k) = m_e \left(\frac{p}{3}\right) \phiG^{-2k}.


    \label{eq:lepton_mass}


\end{equation}


For the electron ($p=3$) we take $k=0$.  Choosing $(p,k)=(5,5)$ yields $m \approx 104.7\,\text{MeV}/c^2$, matching the muon within 1\%.  Similarly, $(p,k)=(11,7)$ gives $m \approx 1579\,\text{MeV}/c^2$, consistent with the tau to within 10\%.  Deviations reflect higher self‑interaction harmonics and are discussed in Appendix~\ref{app:knots}.






\begin{table}[h]


    \centering


    \caption{Comparison of SST lepton masses with experimental values.  Here $p$ is the torus knot crossing number and $k$ the Golden‑layer index.}


    \begin{tabular}{@{}lllll@{}}


        \toprule


        Lepton & Knot $T_{p,2}$ & $k$ & SST mass [MeV] & Physical mass [MeV] \\


        \midrule


        Electron & $T_{3,2}$ & 0 & 0.511 & 0.511 \\


        Muon     & $T_{5,2}$ & 5 & 104.7 & 105.7 \\


        Tau      & $T_{11,2}$ & 7 & 1579  & 1776.9 \\


        \bottomrule


    \end{tabular}


    \label{tab:leptons}


\end{table}






The discrete nature of $k$ reflects the allowed resonant acoustic modes within the Golden‑parameterized NLS core.  Because $k$ increases roughly linearly with the crossing number, heavy leptons appear as topological excitations of the same fundamental trefoil.






%======================================================================


\section{Conclusion}


We have shown that classical hydrodynamics, when endowed with a Golden‑ratio NLS vortex core and const‑druk constraint, yields a self‑consistent description of leptons.  The electron emerges as a degenerate torus whose intrinsic spin gives rise to the Planck constant and whose magnetic moment matches the Bohr magneton with the correct $g$‑factor anomaly.  Discrete scale invariance explains the existence and approximate masses of the muon and tau as topological excitations.  These results suggest that quantum mechanics may be an emergent description of a knotted superfluid, offering a novel bridge between continuum mechanics and particle physics.






\section*{Acknowledgments}


The author thanks the Swirl‑String Theory collaboration for discussions and the developers of open‑source computational tools.  Any errors or omissions remain the author's responsibility.






\appendix






\section{Golden NLS Derivation}


\label{app:golden}


The NLS vortex solution possesses two dimensionless constants $A$ and $B$ that determine the core pressure and translational velocity.  Thermodynamic stability requires that the pressure on the core boundary be uniform and finite.  Setting the net pressure gradient to zero yields $B = A - 1$.  Requiring that the constants be positive and that $B = 1/A$ (for self‑similarity) leads uniquely to $A = \phiG$ and $B = \phiG^{-1}$.  The identity $\phiG - 1 = \phiG^{-1}$ ensures that the velocity and pressure terms satisfy the const‑druk constraint.  This argument underpins the choice of $\phiG$ in Eq.~\eqref{eq:NLS_energy}.






\section{Derivation of Canonical Constants}


\label{app:constants}


Starting from the classical electron radius $r_e = \alpha\hbar/(\me c)$【364253867482672†L183-L189】 and identifying $\rc = r_e/2$, the Hooke closure assumes $m_e c^2 = \tfrac{1}{2} k r_e^2$.  The effective spring constant is therefore $k = 2 m_e c^2/r_e^2$.  The maximal swirl force is $F_{\rm swirl}^{\max} = k\,\rc = m_e c^2/(2\rc)$.  The core density can then be expressed as


\begin{equation}


    \rhocore = \frac{F_{\rm swirl}^{\max}}{\pi \rc^2 \vnorm^2},


\end{equation}


where $\vnorm = \omega_c\,\rc$ and $\omega_c = m_e c^2/\hbar$【620899862539895†L181-L186】.  Evaluating numerically gives $\rhocore \approx 3.893\times 10^{18}\,\text{kg m}^{-3}$.  The effective fluid density at large scales follows by equating the swirl energy density $\rhoe$ with the fluid kinetic energy; solving yields $\rhoF \approx 6.84\times 10^{-7}\,\text{kg m}^{-3}$, consistent with an ultra‑low‑density superfluid.






\section{Torus Knot Mass Spectrum}


\label{app:knots}


The mass spectrum of higher‑order knots can be obtained by considering the line energy per unit length (line tension) of the vortex string and summing over the total length of the knot.  For a $(p,2)$ torus knot, the minimum crossing number is $p$ and the length scales linearly with $p$.  Discrete scale invariance implies that each additional screening layer multiplies the mass by a factor $\phiG^{-2}$.  Combining these effects leads to Eq.~\eqref{eq:lepton_mass}.  A more refined model incorporating the Biot–Savart self‑interaction energy and tangle complexity is discussed in Appendix~C of the original SST manuscript; here we summarise only the leading behaviour.






%======================================================================


\begin{thebibliography}{99}


\bibitem{Hill1894} M.~J.~M. Hill, ``On a spherical vortex,'' \emph{Philosophical Transactions of the Royal Society A}, vol.~185, pp.~213--245, 1894.


\bibitem{Pethick2008} C.~J. Pethick and H.~Smith, \emph{Bose--Einstein Condensation in Dilute Gases}, Cambridge University Press, 2008.


\bibitem{Berloff2004} N.~G. Berloff, ``Interactions of vortices in a trapped Bose--Einstein condensate,'' \emph{Journal of Physics A: Mathematical and General}, vol.~37, pp.~5–19, 2004, \url{https://doi.org/10.1088/0305-4470/37/5/011}.


\bibitem{Conway1970} J.~H. Conway, ``An enumeration of knots and links, and some of their algebraic properties,'' in \emph{Computational Problems in Abstract Algebra}, Pergamon Press, 1970, pp.~329–358.


\bibitem{Kauffman2001} L.~H. Kauffman, \emph{Knots and Physics}, 3rd~ed., World Scientific, 2001.


\bibitem{QEDanom} T.~Kinoshita, ``Theory of the anomalous magnetic moment of the electron,'' \emph{Reports on Progress in Physics}, vol.~59, pp.~1459–1510, 1996.


\end{thebibliography}






\end{document}